\section{Compact estimates and smooth propagation}\label{sec:compact}

We next pass from the stationary estimates to finite-time estimates for the
nonlinear reduced equation.  The essential new point is directional: the
stationary operator is elliptic on the spatial cone containing the projected
trapped covectors, not necessarily over their entire spatial base.  We therefore
retain a full-order bulk estimate for first jets outside that cone.  Combined
with a conic elliptic correction, this gives the same integrable static defect
that is needed in the nonlinear energy comparison. The interaction of
elliptic and integrated wave estimates has a methodological parallel in
Dafermos--Holzegel--Rodnianski--Taylor
\cite{DafermosHolzegelRodnianskiTaylor2024}; the conic estimates for the
reduced Einstein system used here are established in this section.

Throughout this section the solution is smooth on an arbitrary finite bootstrap
slab $[0,T]$.  No constant in an estimate below depends on $T$.  We use the
modulation, targets, reduced equation, stationary family and causal filters
constructed in the preceding section.  In particular,
$u=\partial_t h_{\rm co}$, $\psi$ is the differentiated field in laboratory
components, $p'=a$, $c=\ell_C^p h_{\rm co}$, and
$\bar g=\overline G_{p,c}$ is the instantaneous stationary reference.
All compact pseudodifferential operators act spatially at the current time and
have kernel support strictly inside the computing cylinder.

\subsection{Finite-slab norms and their hierarchy}

For a Banach space $X$, write
\[
 \|V\|_{L^\nu_\alpha X}
   =\|\langle t\rangle^{\alpha/2}V\|_{L^\nu([0,T];X)}.
\]
We may replace $\langle t\rangle$ by $\tau=T_0+t$.  Here $T_0$ is fixed
after all spatial radii, and differences on a fixed initial interval are charged
to the input.  The three levels are
\begin{equation}\label{co2:levels}
 \begin{gathered}
 \begin{array}{c|c|c|c}
 U&U_l&s&A_i\\ \hline
 h_{\rm co}&h_{\rm lab}&10&0\\
 h_{\rm co}&h_{\rm lab}&9&A_H\\
 u&\psi&8&A
 \end{array}\\[3pt]
 A_H=1-\eta_*,\qquad A=1.36,\qquad A_S=1.06.
 \end{gathered}
\end{equation}
where $\eta_*>0$ is sufficiently small.  In particular $2A_H>A$.

Fix a compact physical neighborhood of trapping, disjoint from the horizon,
the inner face, the spatial end and all coefficient probes.  Also choose a
scalar spatial operator
\begin{equation}\label{co2:directional-cutoff}
 B_{\rm nt}=1-\Psi_{\rm tr}
\end{equation}
such that $\Psi_{\rm tr}$ has compact interior kernel support and microsupport
in the stationary-elliptic cone $\mathcal T$.  At high frequency $B_{\rm nt}$
vanishes in a neighborhood of the projections of \emph{both} temporal trapped
sets, with a margin valid for nearby frozen parameters.  Slightly nested
versions will be used without changing notation.  The distinction between
this directional cutoff and the physical trapping neighborhood is important.

For every level in \eqref{co2:levels}, the compact budget $\mathsf K_R$ includes
on fixed enlarged compact sets
\begin{equation}\label{co2:compact-norms}
 \partial_t^bU\in
 L^\infty_{A_i}H^{s-b}\cap L^2_{A_i}H^{s-b-1/2},
 \qquad 0\le b\le s.
\end{equation}
It includes full order $L^2_{A_i}H^{s-b}$ outside the physical trapping
neighborhood, and additionally
\begin{equation}\label{co2:directional-norms}
 B_{\rm nt}\partial_t^bU\in L^2_{A_i}H^{s-b},
 \qquad b=0,1,
\end{equation}
on fixed compact reaches.  Full endpoint order is always retained.  We do
\emph{not} require \eqref{co2:directional-norms} for all higher time jets in the
physical trapping neighborhood.

All compact norm sums use a finite collection of nested reaches and cutoffs.
The collection contains the kernel and probe supports, the error supports in
the estimates below, and the fixed starting overlaps with the far region.
It may be enlarged after the far starting radius $R_0$ is chosen, but is fixed
before the upper matching radius $R$.  Define $\mathsf K_R$ as the sum of the
listed norms on this collection, with fixed positive coefficients, together
with the differential-end norms described next.  Constants on a fixed reach
may depend on it, but not on $R$.  The first-jet norms extend to a
nonincreasing upper cutoff equal to one through $3R$ and zero
after $4R$.  In the differential exterior region they are the energy, bulk and
Hardy norms of
\[
 \widetilde r\,\partial_y^\mu U,\qquad |\mu|\le s-1,
\]
with bulk weight $\widetilde r^{-1-\delta}$ and time weight $\tau^{A_i}$ in
the radial-round measure; $\widetilde r$ is the radius of the constant affine
asymptote of $\bar g$, and $\delta>0$ is arbitrarily small.  We will construct
the corresponding positive energy below.  Mixed norms on other bounded annuli
up to the matching radius may be added with sufficiently small positive
$R$-dependent normalizations.  Their unnormalized estimates are used only in
products or in the choice of $T_0$, never as uniform linearly contracting
controls.  Those normalizations are fixed after $R$ and before the bootstrap
and data smallness.  This convention avoids asserting radius-independent
equivalence of all unscaled mixed norms through the upper shell.

We use a far budget $\mathsf F$ beginning at a fixed large radius $R_0$, chosen
before $R$.  Its complete definition and estimate are given in the next
section; the following components suffice here.  In emitter coordinates
\[
 t=q+r,\qquad z=-\lambda(q)+rn,\qquad |n|=1,
\]
let $T_l=\partial_{t,\rm lab}$, $X=T_l+n\cdot\partial_z$, and
$Z=(rX,\Omega_{q,r})$, with $\Omega_{q,r}$ the sphere rotations.  The good
bulk norms of $(1,Z)Z^I T_l^bU_l$, for $|I|+b\le s-1$ at each
level, in $dt\,dr\,d\omega/r$, have squared weight
\begin{equation}\label{co2:far-weights}
 W_i=r^{p_i-A_i}(t+10r)^{A_i},
 \qquad (p_0,p_H,p_P)=(0.64,1-\eta_*/6,1.62).
\end{equation}
The bad first-derivative norms of $r\partial Z^I T_l^bU_l$ have the weights
of \Cref{fe:F1}; on $t\ge Cr$ the low orders needed in the sliced estimate
are bounded below by $r^{-0.03}\langle t\rangle^{A_i}$.  On this timelike
region we also use
\begin{equation}\label{co2:far-pointwise}
 |h_{\rm lab}|\lesssim\mathsf N r^{-1+0.01}(t/r)^{-A_H/2},
 \qquad
 |\partial h_{\rm lab}|\lesssim\mathsf N r^{-1+0.03}t^{-A_H/2}.
\end{equation}
The relevant co-moving and emitter radii are uniformly comparable.

Set $\mathsf S_a=\|a\|_{L^2_A}$.  Since $A>1$, this controls
$\int|a|$ and hence the change of $p$.  Filtered feedback from the $u$ budget
controls the required positive acceleration derivatives in $L^2_A$, and then
their finite list of endpoint norms.  The total bootstrap size $\mathsf N$
includes $\mathsf K_R,\mathsf F,\mathsf S_a$ and a static budget.
To define it, put $k_{\rm stat}=g_{\rm co}-\overline G_{p,c}$ and let
$K$ run over the fixed compact reaches, with $K_{\rm off}$ the fixed
sets outside the physical trapping neighborhood. Sum the quantities
\[
 \sum_{b=0}^{8}\|\partial_t^bk_{\rm stat}\|_{L^2_{A_S}H^{8-b}(K)}
 +\sum_{b=0}^{8}\|\partial_t^{b+1}k_{\rm stat}
                           \|_{L^2_{A_S}H^{8-b}(K_{\rm off})}.
\]
These are assumed finite-slab budgets at this stage; the sliced estimate
will bound them. All reaches are fixed before $R$. Write $e_{\rm in}$ for the input through the base orders.
Constants denoted $C_{\rm fix}$ may depend on all scales fixed before the
final data smallness, but never on the lifespan.  These are the norm conventions
of \eqref{co2:compact-norms}, with the additional directional control
\eqref{co2:directional-norms}.

The role of the directional control is particularly concrete in the sliced
stationary equation.  On each fixed compact reach its source will have
spatial order six in all directions and order seven after $B_{\rm nt}$.
An order-$-2$ inverse on the complementary
elliptic cone supplies an $H^8$ correction; subtracting it leaves an $H^7$
source for the ordinary one-order stationary estimate.  This recovers an
$H^8$ static residual whose $L^2_{A_S}$ bound, with $A_S>1$, makes its
low coefficient norms integrable in time.  The argument in
\Cref{co2:sliced} uses only the directional first jets specified above.
The low-frequency estimate, proved first, has a different role: it makes
lower compact errors and the linear parameter forcing small enough to absorb
in the final compact energy estimate.

\subsection{The compact low estimate}

Choose a sufficiently large fixed $N_-$ and let $\mathsf l$ be the sum, on
the compact reaches used below, of the $L^2_{A_H}H^{-N_-}$ norms of
$h_{\rm co},\partial_t h_{\rm co}$ and the $L^2_AH^{-N_-}$ norms of
$u,\partial_tu$.  Enlarging $N_-$ merely weakens this norm.

\begin{proposition}[Low-frequency transfer]\label{co2:low-transfer}
Under the simultaneous finite-slab budgets above, there is $\kappa_0>0$ such
that, after choosing $R$ and then $T_0$ sufficiently large,
\begin{equation}\label{co2:low-bound}
 \mathsf l\le C R^{-\kappa_0}(\mathsf K_R+\mathsf F)
       +C_{\rm fix}(e_{\rm in}+\mathsf N^2).
\end{equation}
The exponent can be fixed before enlarging any prior compact reach.  Such
enlargements change constants but not the available radius gain.
\end{proposition}

\begin{proof}
Apply the augmented retarded operator $M_R(\partial_t)$ to $\chi_R U$;
for $U=u$ differentiate the reduced equation first.  The first variation at
the base is exactly $Lh+Q(\partial_t)a$.  After this linear part is removed,
the compact forcing other than $[L,\chi_R]U$ and startup terms has norm
$C_{\rm fix}\mathsf N^2$ in $L^2_{A_i}\overline H^4$; the bar means
extendibility at the outflow boundary.  We spell out why the arbitrary-spin
localization does not alter this fact.  A varying coefficient on a field jet
costs $p,c,h$, their derivatives, or their bounded memories, and never an
uncontrolled position $\lambda$.  At zero field the exact affine-position
freezing identity removes the target defect; its nonlinear remainder therefore
contains a strong acceleration factor.  After one time derivative a changed
coefficient on an undifferentiated-in-time field jet contains $u,\dot c$, or
an acceleration jet.  A product of two weak factors also suffices, since
$2A_H>A$.  There are at most two ordinary differentiations before the four
spatial source derivatives.  These products lie strictly below the half-loss
bulk orders in \eqref{co2:compact-norms}.  The compact kernels preserve the
stated orders, and their distant bounded reaches use the compact-far overlap
norms.  Startup terms and filtered startup tails are bounded by finite-time
stability in terms of $e_{\rm in}$.

Commuting $\tau^{A_i/2}$ through bounded memories and the stable convolution
filter costs $C_R T_0^{-1}$ times the matched budgets, together with initial
data terms.  Apply also the causal smoothing filter
\[
 \left(\frac{c_*}{c_*+\partial_t}\right)^{M_*},
\]
where $M_*$ exceeds the fixed polynomial test losses, including those of a
once differentiated output.  For the shell forcing, the integration by parts
gives the leading Wronskian
\begin{equation}\label{co2:wronskian}
 \int \chi_R' b^{rr}
       (v_*\partial_rU_*-U_*\partial_rv_*)r^2\,dr\,d\omega,
 \qquad r=r_y.
\end{equation}
Here $v_*$ is the transpose test and $U_*$ is the weighted temporal transform.
Replacing both radial derivatives by $\partial_r+\zeta$ leaves the
Wronskian unchanged.  Angular, density and short-range terms have no worse
orders by the Boyer--Lindquist asymptotics; compact additions produce no shell
term.  The transpose radiation estimates and Plancherel thus bound
\eqref{co2:wronskian} by $C c_*^{N'}R^{4\iota}$ times the shell norm of
\[
 (1,\Omega_{\rm co},r_y(\partial_{r_y}+\partial_t))U
\]
in $L^2(\langle t\rangle^{A_i}dt\,dr_y\,d\omega/r_y)$, up to the plateau
errors already listed.

This shell norm is bounded by
\begin{equation}\label{co2:low-shell}
 C R^{-(p_i-A_i)/2}\mathsf F
        +C_{\rm fix}(e_{\rm in}+\mathsf N^2).
\end{equation}
Indeed $y=rn+\lambda(t)-\lambda(q)$, so the two good-direction systems
differ by small parameter coefficients times full first derivatives.  The
component transformation uses $v(t)$, and the difference between $u$ and
the transform of $\psi$ is
$O(v_s-v)\partial h_{\rm lab}+O(a)h_{\rm lab}$.  On this fixed shell the
coefficient differences, also after the required differentiation, have
bounded-memory strong weighted bounds.  Bad derivatives therefore occur only
in products, with allowable $R$-dependent constants.

All tested multipliers are causal: identification on a forward Laplace line
and strip displacement apply separately to the two test kernels produced by
the Wronskian rearrangement.  We can consequently stop forcing at $T$ when
estimating earlier outputs without paying final-time data.  Initial jumps are
absorbed by the smoothing filter.  Taking duals of sufficiently high-order
compact test bounds gives the negative spatial norm defining $\mathsf l$.
Removing the auxiliary filter costs $C c_*^{-1}$ times adjacent local
time-jet budgets and startup transients.  Choose
\[
 c_*=R^{\kappa'},\qquad
 \kappa'N'+4\iota<\frac12\min_i(p_i-A_i),
\]
and then take $T_0$ sufficiently large after $R$.  All $p_i-A_i$ are
positive, so this gives \eqref{co2:low-bound}.  The powers arise only from
the fixed test loss and far weights, not from constants on compact reaches.
Interpolation is consequently available at every order strictly below the
corresponding bulk order.  No full-order bulk estimate on a base-space
stationary-timelike trapping region was used.
\end{proof}

\subsection{Frozen time-local energies}

We now construct the energy supplying the compact budget.  The construction
has two separate roles: a compatible frozen transport gives a half-order loss
at trapping, while strict sign away from both trapped sheets gives the extra
directional first-jet estimate.  Subsequently we will compare the actual
equation with this frozen energy; we will not assume that the actual evolving
metric has a stationary trapped set.

\begin{proposition}[Frozen energy]\label{co2:stationary-energy}
Fix a smooth nearby stationary parameter pair $(p,c)$ and an integer $s\ge8$.
The frozen operator $\mathbf L_{p,c}$ admits a positive time-local first-jet
energy, uniformly for nearby parameters, with full endpoint orders
$H^s\oplus H^{s-1}$.  On each fixed interior reach its favorable derivative
controls
\begin{equation}\label{co2:frozen-dissipation}
 \|\mathcal Y\|_{H^{s-3/2}}^2
       +\|B_{\rm nt}\mathcal Y\|_{H^{s-1}}^2
       -C_s\|\mathcal Y\|_{H^{s-2}({\rm enlarged})}^2,
 \qquad \mathcal Y=(\langle D_y\rangle U,D_tU),
 \quad D_t=-i\partial_t.
\end{equation}
It also controls full mixed physical-nontrapping orders after equation
recovery.  At the inner and outer differential regions the energy is an
ordinary wave current with strict bulk sign and an exiting inner flux.  All
remaining compact derivative errors have total field order at most $s-1$.
The compact additions to the differential currents have strictly interior
spatial kernel support.
\end{proposition}

\begin{proof}
We first construct dissipation near trapping and the horizon, and extend it
through the compact propagation region.  We then join these forms to literal
differential currents near the boundaries.  Finally we recover the original
first jets and record how the energy changes with the stationary parameters.

\paragraph{First-jet reduction.}
With $(p,c)$ fixed, the time-acceleration coefficient is an invertible
multiplication operator plus a compact interior operator of spatial order
$-1$, small at a base order.  Its inverse is multiplication plus negative
order at every smooth order, by the spatial parametrix and derivative gain.
After solving this coefficient and making the elliptic first-jet rescaling
in \eqref{co2:frozen-dissipation}, the evolution has two smooth blocks with
separated scalar real principal roots.  Separation follows from spacelike
slicing, not from timelikeness of the stationary Killing vector.  An
eigenbasis conjugation, followed by the order-by-order off-diagonal equations
dividing only by the root gap, diagonalizes to any required lower order on
compact interior supports.

Write $m_s=2s-2$ for the order of the jet quadratic form.  On each block its
diagonal principal form, expressed on $U$ polarizations, is a positive
homogeneous Hermitian form times the wave action
\[
 |\det\bar g|^{1/2}|\bar g^{t\alpha}k_\alpha|,
\]
of total order $2s$.  Its diagonal energy derivative is ray transport of the
extra Hermitian form, including the nonscalar amplitude transport.  The
scalar geometric-optics divergence is absorbed by wave action and spatial
pushforward.  Equivalently, this is the symmetrized derivative formula in
stationary half-density coordinates for the first-jet system.

Here is the operative compact estimate that this construction will establish.
Let $K$ lie inside a fixed buffered interior reach $K'$, and let $\mathcal Y$
be a smooth first-jet test field supported in $K'$.  Write its normalized
first-jet equation as $\partial_t\mathcal Y=\mathcal G_{p,c}\mathcal Y
+\mathcal F$, where $\mathcal F$ is the full residual, and denote the
constructed positive quadratic form by
$\mathcal E_s=\langle\mathcal A_s\mathcal Y,\mathcal Y\rangle$.
The operator $\mathcal A_s$ has order $m_s$, and
$\mathcal E_s\asymp\|\mathcal Y\|_{H^{s-1}}^2$ on these supports.
For fixed $(p,c)$ the estimate is
\begin{equation}\label{co2:compact-energy-form}
 \begin{split}
 \frac{d}{dt}\mathcal E_s
 &+c_s\|\mathcal Y\|_{H^{s-3/2}(K)}^2
   +c_s\|B_{\rm nt}\mathcal Y\|_{H^{s-1}(K)}^2\\
 &\le C_s\|\mathcal Y\|_{H^{s-2}(K')}^2
       +2\operatorname{Re}\langle\mathcal A_s\mathcal Y,\mathcal F\rangle.
 \end{split}
\end{equation}
No support preservation by the normal-coefficient inverse is assumed:
$\mathcal F$ includes the entire residual.  This formula describes the
buffered compact part; the differential end currents below retain their own
bulk norms and boundary flux.  Higher mixed time jets are subsequently
recovered from the equation and the corresponding source jets.

\paragraph{The half-order loss at trapping.}
Use the compatible fiber metric constructed in the calibrated stationary
family.  Multiply it by
\begin{equation}\label{co2:trap-weight}
 |k(\partial_t)|^{2s-1}b_C,
 \qquad b_C=\exp\bigl(-C(\varphi_+^2-\varphi_-^2)\bigr),
 \qquad \varphi_\pm=z_1\pm z_2.
\end{equation}
Here $C>0$ is a scalar constant, unrelated to the gauge operator.  The normal
variables for the exact Kerr reference satisfy
$V\varphi_\pm=\pm\lambda\varphi_\pm$, with $\lambda>0$.  Exact
transport compatibility therefore gives a negative principal derivative of
relative size
$2C\lambda(\varphi_+^2+\varphi_-^2)$.  Gauge cutoffs below saturation
contribute only fixed lower errors.  This applies to both roots, including
where the spatial base intersects the ergoregion.

The order immediately below this vanishing principal derivative needs care.
In a stationary half-density Weyl chart on a diagonalized block, start the
energy with $e_{m_s}=b_C e^0_{m_s}$ and take its next symbol to be $b_C$
times a $C$-independent choice.  Since the order-one evolution has real scalar
principal symbol $p_{\rm hw}$, its even second Weyl bracket cancels in the
degree-$(m_s-1)$ derivative.  Only first derivatives of $b_C$ can enter from
$e_{m_s}$.  A stationary frame change adds to the next symbol terms linear
in $db_C$, with $C$-independent coefficients.  In this frame the next
symbol has the form
\begin{equation}\label{co2:frame-subprincipal}
 e'_{m_s-1}=b_C r_0+\sum_\alpha(\partial_\alpha b_C)r_\alpha,
\end{equation}
where the normalized symbols $r_0,r_\alpha$ are independent of $C$.
On the trapped set,
\[
 db_C=0,\qquad H_{p_{\rm hw}}(db_C)=0,
\]
the second identity following because $H_{p_{\rm hw}}b_C$ vanishes
quadratically.  More explicitly, for $V=H_{p_{\rm hw}}$,
$V(\partial_jb_C)=\partial_j(Vb_C)-(\partial_jV^k)\partial_kb_C$;
both terms vanish there.  Thus the critical derivative on trapping is bounded
independently of $C$ relative to the unweighted scales.  Off-diagonal errors
are handled by the preceding diagonalization, chosen before $C$.

Factor $C$ from the principal dissipation.  On a normalized symbol patch
let $g$ be the positive Hermitian square root of
$2\lambda b_Ce^0_{m_s}$ and set
$G_\pm=\operatorname{Op}^{\rm w}(g\varphi_\pm)$, chosen self-adjoint.
Both factors have order $m_s/2$ and use the same matrix weight.  Their
squares have principal sum
$2\lambda b_Ce^0_{m_s}(\varphi_+^2+\varphi_-^2)$.
Their degree-$(m_s-1)$ square errors vanish at trapping: at the zero of a
scalar $\varphi$, the first derivatives of $g\varphi$ are $g$ times
scalar derivatives.  For a smooth test section $z$, the elementary inequality
\begin{equation}\label{co2:trap-uncertainty}
 \|G_+z\|^2+\|G_-z\|^2
 \ge \pm\langle i[G_+,G_-]z,z\rangle
\end{equation}
follows from the nonnegativity of
$\|(G_+\mp iG_-)z\|^2$.  On trapping the commutator has, up to the
fixed sign convention, symbol $g^2\{\varphi_+,\varphi_-\}$ of order
$m_s-1$.  The bracket is nonzero with uniform normalized size, and $g^2$
on trapping is positive and independent of $C$.

Retain a fixed fraction of the original $C$-scaled squares and apply
\eqref{co2:trap-uncertainty} to the rest, choosing the sign on each patch.
The resulting positive order-$(m_s-1)$ form dominates the $C$-independent
critical derivative once $C$ is large.  Fix $C$ first and then shrink the
neighborhood on which that form is strict.  Outside it, the retained squares
supply strict principal sign.  Operator-square partitions and stationary
frame changes have on trapping only the critical errors just controlled.
Spatial G\aa rding leaves order $m_s-2$, with constants now allowed to
depend on the fixed $C$.
Consequently the bulk controls $U$ in $H^{s-1/2}$ and $U_t$ in
$H^{s-3/2}$ modulo total field order $s-1$, while the endpoint orders stay
full.  This argument uses symplectic transversality and exact compatible
transport, not stationary ellipticity over a base annulus.

\paragraph{The horizon and compact propagation region.}
At the horizon radial set take action times $b_1^{2s-1}$ and a positive
tensor metric adapted to the horizon-generator null flag.  In generator
parameter, parallel transport scales the null line at rate $-\kappa_+$,
the opposite quotient at rate $\kappa_+$, and rotates the screen quotient.
Flag scaling makes bounded null rotations arbitrarily small.  The rank-two
amplitude norm consequently grows at most at rate $2\kappa_++\epsilon'$.
The local HHV additions are chosen sufficiently small at the center; the
additional high-frequency gauges miss the horizon, and nearby target changes
are small.  The frequency rate is $-\kappa_+$ in the same parametrization,
so
\[
 (2s-1)\kappa_+>4\kappa_+
\]
gives a strict margin, preserved for nearby parameters.  Increasing $s$ only
improves it.

Use the source-neighborhood function and the escape alternatives proved
earlier to extend the principal forms over a large fixed cylinder.  The
remaining compact set, outside smaller trapping and horizon neighborhoods
where the signs have
already been obtained, is covered by finitely many flow tubes.  The ends of
each central bump lie outside this remaining set: they are in exit regions
or in collars with an already strict principal sign.  A stable or unstable
tail may end in the trapping collar; it need not reach a cylinder boundary.
Integrating along the tube gives the required positive derivative, with any
unfavorable endpoint bump confined to those controlled collars.  Their
existing strict signs absorb these endpoint terms after fixed weight choices;
in the horizon collar one uses the increasing source function.
The correction is constant near the exact trapped set, so it does not alter
the preceding critical-order calculation.  Multiplication of the metric by
its decreasing exponential supplies the missing strict signs.  Smooth exact-Kerr models and
finitely many tubes make these choices uniform nearby.  This is a separate
construction on each root block; it does not use positivity of a stationary
Killing current throughout the exterior.

\paragraph{Literal differential currents near the boundaries.}
Outside all compact kernels, the stationary metric has a constant affine
Minkowski leading model $(\widetilde t,\widetilde y)$, with
$\widetilde y=A_{\rm aff}y$ and $\partial_t$ parallel to
$\partial_{\widetilde t}$.  Let $\widetilde r=|\widetilde y|$ and use the
wave stress-energy currents, summed on components, for
\begin{equation}\label{co2:outer-current}
 \chi_\mu=\widetilde r\,\partial_y^\mu U,\quad |\mu|\le s-1,
 \qquad
 N=\partial_{\widetilde t}
       +(f_\infty-\widetilde r^{-\delta})\partial_{\widetilde r},
 \quad 0.9<f_\infty<1.
\end{equation}
The conjugated flat equation has principal part
$-\partial_{\widetilde t}^2+\partial_{\widetilde r}^2
 +\widetilde r^{-2}\Delta_{\mathbb S^2}$.  For
$f=f_\infty-\widetilde r^{-\delta}$, the deformation gives
\[
 f'|\partial_{\widetilde t,\widetilde r}\chi_\mu|^2
 +(2f/\widetilde r-f')\widetilde r^{-2}
                  |\partial_\omega\chi_\mu|^2
\]
in $d\widetilde t\,d\widetilde r\,d\omega$.  Stationary metric deviations
are $O(\widetilde r^{-1+})$, and the first and zeroth order terms in
$\mathbf L$ are respectively $O(\widetilde r^{-2+})$ and
$O(\widetilde r^{-3+})$.  After commuting, the equation recovers
lower-commuted second time jets.  The identity then controls
$\widetilde r^{-1-\delta}|\partial\chi|^2$ with errors
$O(\widetilde r^{-4})|\chi|^2$.  Radial Hardy on each slice controls these
errors and $\widetilde r^{-3-\delta}|\chi|^2$, up to a fixed compact
lower-order anchor.  All first derivatives are controlled, so the affine
tilt of slices is harmless, as is a nonincreasing upper cutoff.

The starting radius $R_c$ for these differential signs can be chosen
independently of $f_\infty\in(0.9,1)$, though final coercivity constants
cannot: $1-f\ge\widetilde r^{-\delta}$ dominates the coefficient
perturbation, while $f'$ and $2f/\widetilde r-f'$ retain their bulk signs.
The coordinates may be chosen from the stationary affine part of the exact
Kerr model, leaving only subleading symbol differences in its radius and
angles.  At the strict spacelike inner outflow use ordinary positive wave
currents with a rapidly increasing radial exponential weight, which supplies
strict damping and the exiting boundary sign.

\paragraph{Matching the currents to the compact forms.}
Far out, use a cutoff in the invariant impact $D/E^2$, equal to one on rays
through the start of the differential region and equal to zero at large
impacts; extend it as one inward.  On its support, sufficiently far out, rays
are close to radial and monotone.  Extend the previous decreasing metric along
them with strict integrable damping $O(\widetilde r^{-1-\delta})$.  Tensor
transport errors relative to the affine model are integrable along such
rays, so this extension remains two-sided comparable to action times
$|E|^{2s-1}$.  The impact bound is fixed from $R_c$ before $f_\infty$ is
taken close to one: in that collar the stationary field is timelike, hence
all relevant impacts are bounded.  Smooth positive-metric transport first
matches the negative derivative in a fixed overlap, and neither its
comparability nor this construction depends on $f_\infty$.

Scale the differential energy by a fixed factor to make it larger than this
extension on incoming rays.  Its characteristic factors relative to stationary
energy approach fixed multiples of $1+f_\infty$ on incoming rays and
$1-f_\infty$ on outgoing rays, with all other jet comparisons uniform.
After that scale is fixed, take $f_\infty$ close enough to one and the
matching radius large enough that the differential energy is smaller on
outgoing rays.  Monotone convex interpolation then has favorable derivative
in both directions.  Toward inner outflow interpolate similarly to a smaller
inner energy.  All these choices precede $R$.

Keep differential currents literally unchanged near each boundary and add
only compact interior quadratic kernels.  Their principal cross terms vanish
in the evolution eigenbasis; the next cross derivative symbol is imposed by
an energy correction dividing by the root difference.  In strict matching
regions, scale cross entries of the differential derivative form by the
products of the interpolation coefficients on its two blocks.  This yields
strict full Hermitian derivative sign there.  Near trapping retain the
diagonalization through the critical order already used.  Spatial kernel
cutoffs lie strictly in the interior and cost smoothing errors on fixed
sets.  The base differential currents may be supported only near their
matching collars, away from trapping and nonlocal PDE kernels.  Their
divergence identities require no acceleration substitution at a boundary;
that substitution is needed only for the additions on buffered interior
supports.  In particular no nonlocal sign on an outflow trace is asserted.

\paragraph{Directional improvement and original jets.}
The strict principal sign off both projected trapped sets allows subtraction
of a sufficiently small positive localized square for
$B_{\rm nt}\mathcal Y$ at order $s-1$.  Its high part vanishes near both
traps with a uniform parameter margin, so this subtraction introduces no
critical symbol there.  One-order composition errors off trapping are
absorbed in the strict principal sign modulo order $m_s-2$; the part
smoothing near trapping costs only the fixed lower orders.  A physical
cutoff avoiding trapping has the same property, and the differential-only
ends already have full estimates.

The order-zero eigenbasis in the first-jet grading and its inverse remain
elliptic even if one temporal root is zero: only their gap has been inverted.
Commuting a scalar spatial cutoff through that stationary eigenbasis costs
one order.  The conic estimate on both diagonal blocks therefore gives
exactly \eqref{co2:frozen-dissipation} on the original first jets.  Very-low
positive coercivity supplements, if required, add only fixed lower-order
derivative errors.  If the stationary parameters now vary, differentiating
this prescribed family contributes at most
$C_R(|a|+|\dot c|)$ times the endpoint norm squared on $r_y\lesssim4R$.
Affine quantities are used algebraically on the frozen wave jets: their
radial coefficient is $\widetilde r$, and no time coordinate factor produces
a spurious $ta$ term.
\end{proof}

\subsection{Sliced control through a stationary-elliptic cone}

The next estimate provides the integrable low coefficient differences needed
for the actual-principal comparison.  Its acceleration estimate must have a
coefficient on $\mathsf F$ independent of $R_0$; otherwise the later compact
and far estimates would form an unabsorbable linear cycle.

\begin{proposition}[Sliced estimates]\label{co2:sliced}
On any fixed enlarged compact reach, the simultaneous budgets imply
\begin{align}
 \mathsf S_a&\le C(R_0)\mathsf K_R+C\mathsf F
                   +C_{\rm fix}(e_{\rm in}+\mathsf N^2),
       &&C\text{ independent of }R_0,\label{co2:sliced-acceleration}\\
 \|g_{\rm co}-\overline G_{p,c}\|_{L^2_{A_S}H^8_{\rm loc}}
   &\le C_{\rm loc}(\mathsf K_R+\mathsf F+\mathsf S_a)
                   +C_{\rm fix}(e_{\rm in}+\mathsf N^2).
                    \label{co2:sliced-static}
\end{align}
The second estimate includes mixed jets of total order eight, and one extra
time derivative with the same mixed order on fixed sets outside the physical
trapping neighborhood.  This extra derivative will supply the mixed
time-radial trace in \eqref{in:trace} when the exterior solution is restricted
to the interior entry cylinder.  Constants on these fixed reaches are
independent of the upper radius $R$.
\end{proposition}

\begin{proof}
We first derive the source regularity, then estimate the acceleration by
adjoint pairing, and finally recover the static field.  At sufficiently late
time choose $\chi_t=1$ for $r_y\le c_1t$ and $\chi_t=0$ for
$r_y\ge2c_1t$, with $c_1>0$ small enough that its support lies in $t\ge Cr_y$.
The fixed earlier interval costs input.  With $p=p(t)$ the sliced equation is
\begin{equation}\label{co2:sliced-equation}
 L_p(0)(\chi_t h_{\rm co})
   =\chi_t(\mathcal R_h-f_B)+[L_p(0),\chi_t]h_{\rm co}.
\end{equation}
The defect $\mathcal R_h$ subtracts the nonstationary reduced equation and
compares all field-dependent terms with their stationary linearized action at
$G_p$.

\paragraph{Compact source orders.}
On a fixed reach the order-four estimate is
\begin{equation}\label{co2:compact-defect}
 \|\mathcal R_h\|_{L^2_A\overline H^4}
 +\|f_B-Q(0)a\|_{L^2_A\overline H^4}
 \le C_{\rm loc}\mathsf K_R+C_{\rm fix}(e_{\rm in}+\mathsf N^2).
\end{equation}
The same source has order six globally, and order seven after $B_{\rm nt}$.
To verify the last, potentially limiting assertion, subtract first at the
frozen target $G_p$ and frozen additions $E_{\rm g}(p,0),C(p,0)$, then
linearize in the field. The linear time-deletion terms have order at
most two and at least one time derivative.
Changed coefficients cost $h,c$, their jets, or acceleration jets and bounded
memories, including $B-G_p$.  True derivatives of the family and its projections
minus frozen derivatives have the same allocation.  Ricci and wave-gauge terms
are quasilinear of order two, with smooth undifferentiated coefficients.
The extra gauges use order-zero or lower spatial kernels and at most one
time derivative apiece.  Thus at most two ordinary field differentiations
precede the spatial source norm, including within compositions.

On the target, the first variation is $Q(\partial_t)a$ with bounded-length
memory.  Its difference from $Q(0)a$ starts with $a'$, controlled by the
$u$ feedback and startup.  The nonlinear difference contains a strong factor
by exact freezing.  A sufficiently large fixed filter order $P$ supplies every
fixed spatial order here.  The coefficient probes are smooth and physically
nontrapped, so their full-order bounds are available.

Order six is below the global half-loss bulk orders: use strong order eight
for $u$, weak order nine for $h$, and $2A_H>A$.  At source order seven the
only limiting factors are
\[
 \partial_y^2 h,\qquad\partial_yu,\qquad\partial_tu.
\]
They have bulk regularity $H^{6.5}$ globally and $H^7$ after $B_{\rm nt}$.
Their coefficient factors are undifferentiated and have endpoint $H^9$.
For a purely spatial second derivative of $h$, its coefficient difference
also has a weak or strong weighted small factor: the exact stationary linear
part has already been subtracted.  We perform this calculation in the reduced
equation \emph{before} inverting its acceleration coefficient.  Hence an
unsmoothed second field derivative has only undifferentiated metric or target
coefficients.  A second-ordinary-derivative placement from an added order-zero
gauge instead carries a time-differentiating kernel of spatial order at most
$-1$; differentiating that kernel in a parameter does not change this spatial
order.  All other connections and inverse-coefficient derivatives distribute
at most first derivatives on individual field factors.  In particular, the
argument never commutes a first-field-derivative coefficient past an unsmoothed
second derivative.

The needed cutoff estimate in three dimensions is
\begin{equation}\label{co2:cutoff-commutator}
 \|[B_{\rm nt},f]v\|_{H^7}
       \le C\|f\|_{H^9}\|v\|_{H^{6.5}}.
\end{equation}
It holds also for matrix-valued $f$, since $B_{\rm nt}$ is scalar.
For completeness, in a compact symbol chart Fourier transform the symbol
difference $b(x,\eta+\ell)-b(x,\eta)$ in $x$.  It decays rapidly in the
dual $x$ variable and is bounded there by a symbol seminorm times
$\min(1,|\ell|/\langle\eta\rangle)$.  After the input Sobolev weight is
removed, the output multiplier is bounded by
\[
 \langle\eta+\ell\rangle^7\langle\eta\rangle^{-6.5}
       \min(1,|\ell|/\langle\eta\rangle)
       \le C\langle\ell\rangle^7.
\]
Indeed, for $|\ell|\le|\eta|/2$ it is at most
$C|\ell|\langle\eta\rangle^{-1/2}$; on the complementary region use
$\langle\eta+\ell\rangle\le C\langle\ell\rangle$.
The smooth $x$-frequency convolution adds only a fixed polynomial weight.
The weighted Fourier
$L^1$ norm of $f$ is bounded by $\|f\|_{H^9}$ because
$9-7>3/2$.  Young's inequality proves \eqref{co2:cutoff-commutator}; smoothing
remainders are harmless.  Commuting the scalar cutoff through a fixed smooth
order-zero matrix kernel similarly gains one order; negative-order kernels
have additional room.  Slightly enlarged supports therefore allow the cutoff
to be placed on the limiting factor in each composition.  Spatially constant
time-dependent coefficients multiply this argument; no time-dependent cutoff
commutator is involved.  This proves the asserted directional source gain.

\paragraph{Far source orders.}
Put $h_l=h_{\rm lab}$ and, on the cut,
\[
 \Delta_a=C\int_{t-C'r_y}^t|a(\sigma)|\,d\sigma
       \lesssim\mathsf S_a t^{-A/2}r_y^{1/2}.
\]
In addition to spatially constant leading combinations of
$T_l\psi$ and $\partial_z\psi$, the far sliced defect is bounded by
\begin{align}\label{co2:sliced-far-source}
 &(r_y^{-1}+\Delta_a)|\partial\psi|+r_y^{-2}|\psi|\notag\\
 &\quad+\Delta_a\bigl(|\partial_y^2h_l|
              +r_y^{-1}|\partial h_l|+r_y^{-2}|h_l|\bigr)\notag\\
 &\quad+|h_l||\partial^2h_l|+|\partial h_l|^2
              +r_y^{-1}|h_l||\partial h_l|+r_y^{-2}|h_l|^2.
\end{align}
This proves the stated source bounds.  To check its applicability,
freeze the affine worldline at the current slice.  Stationary time deletion
then uses $T_l=v(t)\partial_y$, whereas the actual field satisfies
$T_lh_l=\psi+v_s\partial_yh_l$.  The frozen co/lab basis change is spatially
constant.  Comparing in the actual lab equation avoids unsmoothed $v'$ terms.
The first-moment cancellation gives $x-y=O(r_y\Delta_a)$, and the target
jets of order $b\le2$ differ from frozen jets by
$O(\Delta_a r_y^{-1-b})$.  The affine flat leading term cancels exactly;
the Kerr corrections, coefficient derivatives and nonlinear wave-gauge terms
give \eqref{co2:sliced-far-source}.  The cutoff commutator contributes
$O(r_y^{-1}|\partial_yh_l|+r_y^{-2}|h_l|)$ on $r_y\sim c_1t$.

\paragraph{Acceleration pairing.}
Pair \eqref{co2:sliced-equation} with the stationary adjoint kernel
$\mathcal K_p^*$.  The stationary operator term vanishes, and its pairing
with $Q(0)$ has full rank.  Fix a sufficiently large $Y$ before $R_0$.
Replacing $\chi_t f_B$ by $Q(0)a$ on $r_y\le Y$ costs the compact
comparison above.  Beyond $Y$ it costs
$CY^{-1+\epsilon_1}\mathsf S_a$, which is absorbed by the choice of $Y$:
the adjoint kernels are $O(r_y^{-2+\epsilon_1})$,
$\chi_t f_B=O(r_y^{-2}\chi_t\widetilde a)$, and enlarged averaging over
comparable radii is uniformly bounded in the strong weighted norm on the cut.
Choose $0<\epsilon_1<\zeta_1<0.13$, with $\zeta_1$ small.

For leading terms beyond $CR_0$, write $T_l=X-n\cdot\partial_z$ and
integrate the remaining spatial derivative once on the time slice, as in
\Cref{fs:uniform-coupling}.  The test coefficient is
$O(r_y^{-2+\epsilon_1})$ and gains an inverse radius upon differentiation.
This includes derivatives of the cutoffs and $n$, since the emitter-coordinate
first derivatives require only bounded velocities.  It therefore suffices to
control
\[
 \|r^{\zeta_1}t^{A/2}(rX\psi,\psi)\|_{L^2(dt\,dr\,d\omega/r)}
       \le C\mathsf F.
\]
Here $C$ is independent of $R_0$, precisely because $p_P-A=0.26$.

For the other far terms $S'$ in \eqref{co2:sliced-far-source}, use directly
the norm of $r_y^{1+\zeta_1}t^{A/2}S'$ in this measure.  The bad norm of
$r_y\partial\psi$ loses at most $0.015$ powers of radius, compensated by
$r_y^{-1}+\Delta_a$.  Using tangential $r^{-1}Z$, $X$, and
$T_lh_l=\psi+v_s\partial_yh_l$, with $v_s$ small, gives
\[
 |\partial_y^2h_l|
 \lesssim |\partial\psi|+r_y^{-1}|Z^{\le1}\psi|
                         +r_y^{-2}|Z^{\le2}h_l|.
\]
The differentiated coordinate coefficients are $O(r_y^{-1})$, including
$\partial v_s$.  In weak terms use
$t^{A/2}\Delta_a\lesssim\mathsf S_a r_y^{1/2}$.  For quadratic terms
combine \eqref{co2:far-pointwise} with weak bulk norms of
$h_l,r_y\partial^{1,2}h_l$, losing at most $0.015$ powers.  For example the
extra radius factor for $h_l\partial^2h_l$ is bounded by
\[
 C\mathsf N r_y^{\zeta_1+0.015-1+0.01+(A-A_H)/2},
\]
a negative power.  The cutoff commutator has the required extra inverse
radius.  This proves \eqref{co2:sliced-acceleration}; dependence on $R_0$
occurs only in the fixed compact and starting regions.

\paragraph{The conic elliptic correction.}
Apply stationary Fredholm regularity at growth power $\gamma=0.96$ to
$\chi_t h_{\rm co}-H_C(p)c$, with the seven decaying-mode measurements and
the ten constant-asymptote measurements from the preceding section.  The
remaining measurement costs
$C_{\rm loc}(\mathsf S_a+\mathsf K_R+e_{\rm in})$: only $\ell h_{\rm co}$
remains, controlled because $F(0)$ is invertible and
$(F(\zeta)-F(0))/\zeta$ is stable on
$\partial_t(\chi_{\rm st}\ell h_{\rm co})$.

At infinity the source needs only
$r_y^{2-\gamma}t^{A_S/2}S'$ in $L^2(dt\,dr_y\,d\omega/r_y)$, with
equation derivative order zero.  The direct bounds just proved supply it,
including $\chi_t f_B$ by strong averaging.  Even the unsubtracted leading
terms fit, since
\[
 1-\gamma+0.015<(A-A_S)/2.
\]
On a fixed enlarged compact reach let $f$ denote the full sliced source.
We have established $f\in H^6$ and $B_{\rm nt}f\in H^7$, with the
claimed time-weighted norms.  Choose a properly supported order-$-2$
parametrix $\mathcal Q_p$ for $L_p(0)$ on an enlarged microsupport of
$\Psi_{\rm tr}$, so that
\[
 L_p(0)\mathcal Q_p\Psi_{\rm tr}
       =\Psi_{\rm tr}+\mathcal R_p,
\]
where $\mathcal R_p$ is smoothing to any prescribed finite order.  Include
the compact pseudodifferential lower terms in the full symbol construction;
finite-rank terms and off-cone leakage are smoothing.  The correction
$\mathcal Q_p\Psi_{\rm tr}f$ is $H^8$, and its fixed coefficient
measurements have the same bound.  Subtracting it leaves the $H^7$ source
$B_{\rm nt}f-\mathcal R_pf$.  The ordinary stationary estimate with its
one-order gain now gives $H^8$ for the remainder.

The stationary application therefore uses two spatial regularity regimes:
source order six with the additional $B_{\rm nt}$ order seven on the fixed
compact reaches, and only the weighted source order zero specified above
near infinity.  Switch between them across a fixed exterior annulus where
the \emph{full} spatial symbol is elliptic, retaining a margin of high source
regularity through the switch.  The stationary propagation and radial
estimates inside, and $b$-ellipticity with the same indicial gap outside, are exactly the
variable-order observation in \Cref{st:stationary-estimates}.  Lowering far
derivative orders introduces no new kernel ambiguity.  This argument needs
conic ellipticity near projected trapping and ordinary ellipticity far out;
it does not assume ellipticity on an entire trapping base annulus.

Finally, replacing $G_p+H_C(p)c$ by $\overline G_{p,c}$ costs quadratic
$c$ terms; replacing $B$ by $G_p$ costs bounded acceleration memories.
Mixed jets with a time derivative are obtained directly from $u$, not by
differentiating the parametrix: they use order seven globally and full order
eight outside physical trapping.  Positive derivatives of the reference and
projections use strong parameter and field jets through order nine for $c$.
At order nine the leading term is $\ell_C^p\partial_t^8u$ at the physically
nontrapped probes; the highest strong order is needed only in bulk.  This
proves \eqref{co2:sliced-static} and its extra-time-derivative assertion.
\end{proof}

\subsection{Comparison with the actual principal part}

The frozen energy alone does not estimate the nonlinear equation.  Its
principal wave coefficients must be replaced by the actual ones without
losing a derivative.  The sliced estimate makes the lower coefficient
differences integrable; the differential principal difference is instead
absorbed algebraically into the energy.

\begin{proposition}[Actual-principal comparison]\label{co2:actual-comparison}
For each level $(U,s,A_i)$ in \eqref{co2:levels}, let
$d_M^{\alpha\beta}$ be the change from $\bar g$ to $g$ of the differential
order-two wave coefficient in stationary normalization.  On fixed compact
reaches the reduced equation has the form
\begin{equation}\label{co2:source-ledger}
 \mathbf L_{p(t),c(t)}U+d_M^{\alpha\beta}\partial_{\alpha\beta}U
       =S_i^{\rm lin}+S_i^{\rm err},\qquad
 S_i^{\rm lin}=\pm Q(\partial_t)b_i,\qquad
 b_i=a,a,a',
\end{equation}
up to common normalizations, where
\begin{equation}\label{co2:error-source-space}
 S_i^{\rm err}\in L^1_{A_i}H^{s-1}+L^2_{A_i}H^{s-1/2},
 \qquad
 \|S_i^{\rm err}\|\le C_{\rm fix}\mathsf N^2
                         +C_{\rm fix}e_{\rm in}.
\end{equation}
On larger bounded differential annuli only $L^2_{A_i}H^{s-1}$ is required.
There is an actual-principal correction of the frozen energy whose additional
integrated derivative errors are bounded in squared estimates by
$C_{\rm fix}\mathsf N^3$, up to input terms.  The frozen half-loss bulk and
directional improvement remain on the left side of that estimate.
\end{proposition}

\begin{proof}
We compare the frozen and evolving equations, retaining the actual
principal coefficients and checking every source derivative count.

\paragraph{Integrable low differences.}
Acceleration and $\dot c$ have strong $A$-weighted bulk and endpoint bounds
at the required differentiated orders; $c$ has weak $A_H$ bounds and
unweighted high bounds.  In detail, derivatives of $c$ of orders one through
nine use the full strong $u$ jets at the probes.  The unweighted top $h$
control also supplies $c^{(10)}$, but no strong weighted estimate for this
jet is used.  The mixed recovery below differentiates the $h$ equation at
most eight times and the differentiated $u$ equation at most six times;
its parameter placements require at most $c^{(9)}$.  Derivatives of $p$ use
the fixed smooth filter and startup bounds.  On compact sets $B-G_p$ and true derivatives
minus affine-position frozen derivatives are bounded strong acceleration
memories, with no position growth.  Put
\begin{equation}\label{co2:recentered-fields}
 \xi=\bar g-G_p,\qquad k_h=h_{\rm co}-\xi,\qquad X_0=B+\xi.
\end{equation}
The reference increment $\xi=O(c)$ has fixed smooth spatial orders.
By Proposition~\ref{co2:sliced}, $k_h$ and $g_{\rm co}-\bar g$ have
integrable low norms of size $C_{\rm fix}\mathsf N$, apart from initial
terms.  Indeed $A_S>1$ and Cauchy--Schwarz turn their weighted $L^2$ mixed
order-eight bounds into $L^1_t$ bounds.  These control spatial Lipschitz norms
and the first true time derivative of low principal differences.  Likewise
$|a|+|\dot c|$ and the true-minus-frozen reference derivatives are integrable.
High placements are controlled by their ordinary field and coefficient
budgets, not by an asserted integrable high norm.  On larger bounded annuli
the full physical-nontrapping and far overlap controls suffice with
$R$-dependent constants in products, including the strong comparison of $u$
and $\psi$.

\paragraph{The nonlinear source ledger.}
For $U=h_{\rm co}$, subtract the reduced expression at $g=B$, namely $f_B$.
Treat target, gauge additions and compensator as fixed coefficients and
inputs, also within kernel applications.  Taylor-expand both
$X_0+k_h$ and $X_0-\xi$ at the \emph{same} $X_0$.  For the reduced
differential and kernel expression $\mathcal D_{\rm red}$ the field-dependent
part is
\begin{equation}\label{co2:taylor-recentering}
 \mathcal D_{\rm red}'(X_0)h_{\rm co}
       +\operatorname{Rem}_{X_0}(k_h)
       -\operatorname{Rem}_{X_0}(-\xi).
\end{equation}
Each remainder is at least quadratic in its displayed increment and jets,
with smooth coefficient multiplicities.  This common expansion point avoids
an unabsorbed high first jet multiplied only by a single low slowly decaying
$\xi$.  By stationary tangency, the low residual gauge of $X_0$ at freezing
is $O(c^2)$; away from freezing it also has integrable acceleration memories
and parameter derivatives, including the compensator.  The linear term in
\eqref{co2:taylor-recentering} is $\mathbf L_{p,c}$, up to the differential
principal change and smooth integrable lower coefficient errors.  In
particular the omitted derivative of the wave-gauge operator acting on the
gauge of $X_0$ is of this class, since $H_C$ solves the stationary linear
gauge equation.

The smooth $-\xi$ remainder is quadratic.  In the other remainder count the
changed leading coefficient on $\partial^2k_h$ in
$d_M\partial^2h_{\rm co}$; its leftover action on $\partial^2\xi$ needs
at most order-zero $h$ factors.  Remaining ordinary products have at most
one derivative on each $k_h$ factor, or a high undifferentiated field times
fixed smooth reference derivatives.  A low $k_h$ factor through order two,
with a fixed Sobolev buffer, is integrable.  Thus high first derivatives fit
$L^1_{A_i}H^{s-1}$ using endpoints, and high undifferentiated factors fit the
half-loss bulk-source space.  Smooth quadratic $\xi$ products use
$2A_H>A$.

For $U=u$, differentiate the full $h$ equation.  Its field derivative is
at $g$, with the principal difference as in \eqref{co2:source-ledger}.
Lower coefficient differences, including the residual-gauge terms, remain
integrable at low order.  A high order-two $h$ jet multiplied by $u$ or a
differentiated coefficient drift uses unweighted $H^{7.5}$ bulk from the top
mixed budget and a strong low endpoint.  A high first derivative of $u$
uses an integrable low difference.  A high differentiated $h$ factor and a
low differentiated $u$ factor multiply their unweighted and strong $L^2_t$
budgets into $L^1_{A_i}H^{s-1}$.  Derivatives of $B,E_{\rm g},C,\Theta$
relative to the zero-field expression contain strong coefficient factors.
The spatial tame product estimates in dimension three therefore give the
claimed source orders.  The following representative borderline placements
record the distinct uses of endpoints and bulk norms.  Each low factor uses
the fixed spatial multiplication buffer in its compact budget; integrable
coefficient differences use the sliced bound.
Here $K_{-1}$ denotes a kernel of spatial order $-1$, and $d_1$ an
undifferentiated coefficient difference.
\begin{center}
\small
\begin{tabular}{@{}p{0.21\linewidth}p{0.43\linewidth}p{0.25\linewidth}@{}}
\toprule
Placement & Factors used & Source space \\
\midrule
High first derivative of $k_h$ times a low derivative, at the $h$ levels
& Weighted endpoint $H^{s-1}$ times a low $L^1_t$ factor
& $L^1_{A_i}H^{s-1}$ \\[3pt]
$u\,\partial^2h$ at $s=8$
& Unweighted $H^{7.5}$ bulk for $\partial^2h$ times a strong low endpoint
& $L^2_AH^{7.5}$ \\[3pt]
$K_{-1}(d_1\partial^2u)$ at $s=8$
& Input in $L^1_AH^6$: low integrable $d_1$ times a weighted endpoint,
  or unweighted high bulk $d_1$ times strong low bulk
& $L^1_AH^7$ after the spatial gain \\
\bottomrule
\end{tabular}
\end{center}
Opposite high/low placements and smooth coefficient multiplicities obey the same tame product estimates.

The nonlocal terms use the spatial gain in the last row.  The added
operators have spatial orders zero and $-1$; a time hit carries order $-1$.
They occur in compositions of at most two first-order gauge factors with
smooth wave coefficients.  Any changed coefficient on a remaining order-two
field term beyond the differential principal change therefore has a spatial
gain.  For example, an order-$-1$ kernel on $d_1\partial^2u$, where $d_1$
is an undifferentiated coefficient difference, uses either integrable low
$d_1$ and the weighted endpoint $H^{s-2}$ norm of $\partial^2u$, or
unweighted bulk $H^{s-2}$ for $d_1$ and strong low bulk for
$\partial^2u$.  The output is $L^1_{A_i}H^{s-1}$.  For $h$ split
$k_h+\xi$ and use half-loss bulk for a high coefficient multiplying a
slowly varying spatially smooth factor.  The same argument applies when the
negative-order kernel precedes multiplication, using a fixed low Sobolev
buffer, or appears on both gauge factors.  Hits on kernel parameters use the
weighted parameter derivatives already counted.  Finally, after its linear
part is removed, $f_B$ and $\partial_t f_B$ contain a strong acceleration
factor, jet or memory on $r_y\le4R$, by exact affine-position freezing.
This proves \eqref{co2:error-source-space}.

\paragraph{The separate linear source.}
It must not be treated as a nonlinear error.  On prior fixed sets its
required $L^2_{A_i}$ spatial orders are bounded by
$C\mathsf l+C_{\rm fix}e_{\rm in}$ through filtered feedback, from
$\ell h,\ell h,\ell u$ respectively.  Enough fixed filter gain also gives
the required endpoint orders.  The differential exterior estimate uses
\begin{equation}\label{co2:outer-linear-source}
 \left(\int^{O(R)}\widetilde r^{1+\delta}
  \|\widetilde r\,\partial^{\le s-1}Q(\partial_t)b_i
       \|_{L^2_{A_i}L^2_\omega}^{2}\,d\widetilde r\right)^{1/2}
 \lesssim C R^{\delta/2}\mathsf l+C_{\rm fix}e_{\rm in}.
\end{equation}
Explicitly, $Q$ and its required derivatives cost
$r_y^{-2}$ times bounded averaging of filtered jets.  Taking $T_0\gg R$
makes the plateau weights comparable across those memories; slowly varying
affine measures are uniformly comparable as well.

\paragraph{Absorption of the principal difference.}
Use the actual wave metric in the differential currents.  On fixed compact
sets the principal commutator differences obey the tame split
\begin{equation}\label{co2:principal-commutator}
 C\|\nabla_y d_M\|_\infty\|\partial^2U\|_{H^{s-2}}
       +C\|d_M\|_{H^{s-1}}\|\partial^2U\|_{\rm low}.
\end{equation}
The first term is integrable low times endpoint.  For a high coefficient
with $h$, split $k_h+\xi$ as above, allowing the extra half order on the
smooth $\partial^2\xi$ placement; with $u$, multiply unweighted and strong
bulk norms.  Derivatives of the principal-current coefficients cost current
order times integrable low.  On larger bounded annuli use full mixed bulk at
$R$-dependent constants.  The strict inner outflow remains exiting, including
when a cutoff differential current by itself has no sign in the interior.

On buffered compact kernel supports, the solved first-jet evolution is the
complete frozen evolution plus $iT_{\Delta P}\mathcal Y$, modulo the source
classes already listed.  Here $T_{\Delta P}$ is the spatial paradifferential
order-one change of the differential principal symbol.  To see this, first
normalize actual and frozen equations by their differential time coefficients.
Their inverse difference times stationary lower terms belongs to the same
source classes; the extra stationary acceleration term has spatial gain.
Inverting its normalized $I+P_{-1}$ contributes only an order-$-1$
correction on the differential principal difference.  The non-paralow
remainders cost high coefficient $H^{s-1}$ times low $\partial^{\le2}U$,
with the smooth $\xi$ placement permitted at order $s-1/2$.  These are the
preceding product estimates.  Stationary cutoff errors stay within the
stationary calculation, so no nonlocal side condition is introduced.

Conjugate the first jets by $\langle D_y\rangle^{s-1}$ on these supports
and put $z=\langle D_y\rangle^{s-1}\mathcal Y$.  Write the actual
differential evolution symbol as
\[
 p_d=S_d\operatorname{diag}(\lambda_+\mathrm{Id},\lambda_-\mathrm{Id})S_d^{-1},
\]
where $S_d$ depends algebraically and smoothly on the wave coefficients.
Let $e_\pm$ be the stationary diagonal principal forms for the compact
addition to the base differential currents; they need not be positive
separately.  Correct this addition at spatial paradifferential principal
order by
\begin{equation}\label{co2:algebraic-pullback}
 A_d=S_d^{-*}\operatorname{diag}(e_+,e_-)S_d^{-1}.
\end{equation}
Keep the lower stationary prescriptions.  Since the roots are real scalar
roots, the algebraic identity
\begin{equation}\label{co2:principal-symmetrization}
 A_dp_d=p_d^*A_d
\end{equation}
cancels the order-one part of $i(A_dP_d-P_d^*A_d)$ exactly.  The
differential wave currents separately symmetrize their actual and frozen
symbols, so the same cancellation holds for the total corrected energy.

After this exact cancellation, only one-order symbolic accuracy is needed.
For the symbolic calculus, see Bony
\cite[\S 3, Theorems~3.2--3.3]{Bony1981}; a quantitative formulation
of the one-spatial-Lipschitz-derivative composition and adjoint bounds is
\cite[Definition~2.1 and Theorem~2.6(ii)--(iii), $\rho=1$]{AlazardBurqZuily2014}.
The first-order spatial Bony composition and adjoint estimates
\cite[Theorems~2.4.11--2.4.12]{Metivier2003} leave an order-zero
difference from the stationary derivative.  After the source
pairings already estimated, its quadratic form is bounded by
\begin{equation}\label{co2:integrable-energy-error}
 C_s\bigl(\|d_M\|_{W^{1,\infty}}+\|\partial_td_M\|_\infty
                  +|a|+|\dot c|\bigr)\|z\|_{L^2}^2.
\end{equation}
The fixed low regularity buffer from mixed $H^8$ control supplies the
spatial coefficient and smoothing-remainder bounds in these rules.
Normalized frequency derivatives use finitely many smooth angular symbol
seminorms, with the same low coefficient bounds; their number does not grow
with the energy order.  The algebraic pullback has one true time derivative,
acting on parameter values and low variation.  Stationary subprincipal terms
and compact kernel-cutoff interactions with the paradifferential change have
at most current order.  The coefficients of $b_C$ have already been fixed.
Smallness preserves endpoint positivity, and integration of
\eqref{co2:integrable-energy-error} against $\tau^{A_i}$ costs
$C_{\rm fix}\mathsf N^3$ in squared estimates, up to input.

Thus the corrected energy retains the frozen dissipation, including the
$B_{\rm nt}$ first-jet term in \eqref{co2:frozen-dissipation}, with only an
integrable endpoint error.  It requires neither a trapping sign for the
actual metric nor full directional estimates on higher mixed time jets.
The scalar cutoff has not been commuted through the changing order-two wave
equation, and every compact kernel remains inside its buffered support,
away from the outflow trace.
\end{proof}

\subsection{The integrated compact estimate}

\begin{proposition}[Compact high estimate]\label{co2:compact-high}
Under the simultaneous finite-slab budgets, for some $\kappa>0$,
\begin{equation}\label{co2:high-bound}
 \mathsf K_R\le C R^{-\kappa}(\mathsf F+\mathsf K_R)
       +C_{\rm fix}\bigl(e_{\rm in}+\mathsf N^{3/2}
                                 +(e_{\rm in}\mathsf N)^{1/2}\bigr).
\end{equation}
The acceleration and sliced static bounds enter this estimate only as
coefficient smallness and nonlinear errors, not as additional linear forcing
terms.  The choices are the stationary energy parameters, then a sufficiently
large $R$, then the plateau $T_0$, and finally the bootstrap and data
smallness.
\end{proposition}

\begin{proof}
Apply Proposition~\ref{co2:actual-comparison} to the frozen energies at the
three levels \eqref{co2:levels}.  Put a nonincreasing cutoff, with derivative
$O(R^{-1})$, on the outer differential current. Its unfavorable squared
flux is bounded, up to finite-time costs, by
\begin{equation}\label{co2:high-shell-flux}
 C\sum_{|\mu|\le s-1}
 \|\tau^{A_i/2}(1,\Omega,rX)\partial_y^\mu U
         \|^2_{L^2(dt\,dr\,d\omega/r;\,\mathrm{shell})}.
\end{equation}
On a pure outgoing flat covector the flux exits strictly because
$f_\infty<1$.  After the matching currents are fixed, the affine, cutoff
normal and cone changes are small for large $R$ and a small bootstrap.
Finite-dimensional definiteness therefore bounds every unfavorable component
by the good derivatives of $\chi_\mu$; differentiating its radius factor
accounts for the order-zero term in \eqref{co2:high-shell-flux}.

The square root of \eqref{co2:high-shell-flux} is at most
\begin{equation}\label{co2:high-shell-bound}
 C R^{-\kappa_b}\mathsf F+C_{\rm fix}(e_{\rm in}+\mathsf N^2),
 \qquad
 0<\kappa_b<\min_i\bigl\{(p_i-A_i)/2,0.98\bigr\}.
\end{equation}
Indeed
\[
 \partial_{z_d}=-n_dT_l+r^{-1}(n_drX+\Omega_d)
\]
expands a spatial derivative string into
$r^{-b}c(n)Z^IT_l^j$, with $j+|I|\le|\mu|$ and $b\ge|I|$ at the flat
level.  Order the lab time derivatives before $Z$ derivatives.  Perturbed
emitter and co-frame coefficients add only small coefficient products
controlled by the fixed-annulus overlap norms.  The $I=0$ terms use the
good-weight gain in \eqref{co2:far-weights}; every $I\ne0$ term has an
inverse radius, and the bad weights
$\gtrsim t^{A_i}r^{-0.04}$ then suffice.  At the strong level the difference
between $u$ and the transformed $\psi$ contains strong bounded memories
such as $v(t)-v_s,a$, and at most nine derivatives of $h_{\rm lab}$.
At the $h$ levels any time derivative of the component change carries strong
acceleration and lower fields.  Thus full good and bad first-jet far controls
through $j+|I|\le s-1$ suffice.  The earlier fixed interval, allowed to
depend on $R,T_0$, costs input by finite-order Cauchy stability.

All fixed compact lower derivative errors have total spatial or first-time
order at most $s-1$, strictly below the half-loss bulk.  At matched time
weights interpolate them by
\begin{equation}\label{co2:lower-interpolation}
 R^{-\iota_0}\mathsf K_R+C R^{K\iota_0}\mathsf l,
 \qquad K\iota_0<\kappa_0/2,
\end{equation}
where $K$ is a sufficiently large fixed interpolation constant and
$\kappa_0$ is from Proposition~\ref{co2:low-transfer}.  Very-low-order
coercivity supplements fit the same estimate.  After substituting
\eqref{co2:low-bound}, the data and quadratic pieces have only
$C_{\rm fix}$ costs.  The linear forcing pairing from
\eqref{co2:outer-linear-source} is at most
\[
 C R^{\delta/2}\mathsf l\,\mathsf K_R
       +C_{\rm fix}e_{\rm in}\mathsf N.
\]
Choose $0<\delta<\kappa_0$ sufficiently small.  The low-transfer exponent
is already fixed before these enlarged compact reaches: their constants and
the outer current's slack power do not worsen its test loss or far-weight
gain.

Differentiating the time weight uses the unweighted top bulk for the extra
half derivative.  On a fixed compact set, for example,
\begin{equation}\label{co2:plateau-interpolation}
 \tau^{A_i-1}\|U\|_{H^s}^2
 \lesssim\tau^{A_i/2-1}
       \bigl(\tau^{A_i/2}\|U\|_{H^{s-1/2}}\bigr)
       \|U\|_{H^{s+1/2}}.
\end{equation}
The analogous first-time-jet estimate shifts the spatial order by one.
For the strong level it uses $u_t=\partial_t^2h$ in unweighted
$H^{7.5}$ bulk.  On farther bounded annuli the permitted radius-dependent
constants and mixed top orders suffice.  Since $A_H,A<2$, taking $T_0$
large after $R$ makes these coefficients as small as required.  Integration
up to arbitrary endpoints, using \eqref{co2:high-shell-bound},
\eqref{co2:lower-interpolation} and the cubic comparison errors, proves
\eqref{co2:high-bound} first for spatial and first-time jets, with the full
mixed budget still on its right.

It remains to recover that mixed budget.  Successively solve the equations
with the stationary normal coefficient inverse, including its order-$-1$
addition, after $j$ time derivatives, for $0\le j\le s-2$.
Lower time jets in the stationary homogeneous part have already been bounded;
parameter derivatives on coefficients are products or input.  A nonlinear
term, including a principal difference on the current time jet, needs only
bulk order
\[
 r'=s-j-2-\theta,
 \qquad\theta=1/2\text{ globally},\quad
 \theta=0\text{ in a full-order region},
\]
and endpoint order $s-j-2$.  Its norm is
$C_{\rm fix}\mathsf N^2$ by the bootstrapped mixed estimates.  Explicitly,
if $k$ time derivatives hit an undifferentiated field coefficient on an
order-two jet, put that coefficient in endpoint $H^{s-k}$ and the other
factor in bulk $H^{r'+k}$.  The product estimates have the margins
\[
 r'\ge-1/2,\qquad s-k\ge r'+2,\qquad
 (s-k)+(r'+k)-r'=s.
\]
Two first-order placements and further smooth coefficient factors have the
same margin.  Place the matched time weight on one factor and use two
endpoints for endpoint output.  Products containing $u$ use unweighted top
$h$ and strong $u$; drifts involving $h$ alone carry strong differentiated
parameters.  Kernels preserve or improve spatial order and have at most two
time derivatives before this mixed differentiation.  The preceding parameter
counts cover all cases.  The linear source is still filtered low measurement
with fixed compact costs.

For full mixed bulk outside physical trapping, use nested smooth spatial
cutoffs with margins.  Place the cutoff on the bulk factor in products;
commuting it through stationary kernels gains a full spatial order, more
than the global half-order loss.  No full $B_{\rm nt}$-cut estimate for
$j\ge2$ is needed inside the remaining physical neighborhood.  Additional
mixed norms on larger annuli follow at $R$-dependent cost and are included
with the small normalizations specified at the start; they occur elsewhere
only in products or plateau choices.  Thus the left side becomes the full
$\mathsf K_R$, proving \eqref{co2:high-bound}.
\end{proof}

\subsection{Higher orders after base closure}

The next statement is used only after the compact and far estimates have been
closed at the fixed base orders in the next section.  It is important that
it impose no smallness on higher data seminorms.

\begin{proposition}[Subexponential smooth bounds]\label{co2:higher}
Assume the global base bounds resulting from the compact-far closure.  For
every fixed $C_0>0$, every fixed smooth order, and every $\sigma>0$, all
metric, coefficient and parameter norms at that order on
\[
 0\le t\le T',\qquad r_y\le C_0(1+T')
\]
are bounded by $C_\sigma e^{\sigma T'}$.  The constant may depend on the
smooth datum, the order and $C_0$, but the base neighborhood and fixed filter
do not.  We denote these bounds by $e^{o(T')}$.
\end{proposition}

\begin{proof}
We use induction in derivative order, with the conic energy and its
parameter counts explicit.  The closed base estimates imply, uniformly on
the slice in the low unscaled supremum norms needed below, through at least
two field derivatives,
\begin{equation}\label{co2:low-decay}
 h=o(1),\qquad g_{\rm co}-\bar g=o(1),\qquad
 c,\dot c,a=o(1),\qquad t\longrightarrow\infty.
\end{equation}
On fixed compact sets use weighted endpoints; at large radius use the far
point estimates, including \Cref{fe:F7} and \Cref{fe:F8}, with their spare
rotations.  They give uniform decay of low unscaled lab derivatives as
$r\to\infty$, while the co transformation adds bounded velocities and
differentiated parameters.  Fixed annuli have decay in $t$.  Also
$B-G_p=o(1)$: compactly this is a bounded memory, and near infinity both
leading co asymptotes are exactly the affine flat metric $\eta_{v(t)}$,
with symbol-decaying differences and bounded differentiated frame changes.
True low reference derivatives minus their formal frozen counterparts are
$o(1)$.  The affine changes and symbol remainders in $\bar g-G_p$ use
$c=o(1)$.  No convergence of the coordinates is required, although $p$
does converge because $a\in L^1$.

Fix a terminal time $T'$ and work on shrinking cylinders with outer radius
\begin{equation}\label{co2:higher-cylinder}
 C_0(1+T')+C_1(T'-t),
\end{equation}
where $C_1$ is large enough for strict exit and all compact kernels have a
buffer inside.  Induct on total order, using a larger such cylinder for the
previous order if needed.  Orders through seven start the induction: bulk
bounds through order eight convert to unscaled $L^2$ bounds with at most
polynomial losses in $T'$ from the measures and radial weights on these
growing balls.  Conversion of derivatives uses lab time, $r^{-1}Z$, bounded
velocities and the available filtered acceleration jets.  Enlarged spatial
regions and backward unit slabs give order-seven endpoints; startup uses
fixed finite-time bounds.  The spare rotations in the far point estimates
ensure the uniform low decay used in \eqref{co2:low-decay}.

At each $s\ge8$, Proposition~\ref{co2:stationary-energy} is available near
the limiting stationary parameters.  Exact compatible saturation at trapping
is independent of $s$, and the horizon inequality becomes stronger.  For
this one order, start at a sufficiently late time that the actual-principal
correction preserves positivity; the preceding fixed interval has smooth
symbol-tail estimates by finite-time propagation.  Continue the differential
current to the exiting boundary \eqref{co2:higher-cylinder}, with no
matching-shell loss.  Its top-order energy equivalence constants are
independent of the terminal radius, since
\begin{equation}\label{co2:affine-density}
 \|\widetilde r G\|_{L^2(d\widetilde r\,d\omega)}^2
       =|\det A_{\rm aff}|\|G\|_{L^2(d^3y)}^2,
\end{equation}
and $f_\infty<1$ is now fixed.  Lower terms from differentiating
$\chi_\mu$ use the preceding induction.  The parameter derivatives of
$\widetilde r$ relative to $\widetilde r$, of the affine angles, and of
the density ratio are uniformly bounded, so differentiating current
coefficients incurs no terminal-radius factor.  The frozen derivative is
favorable modulo the fixed compact lower order already estimated.

In the principal comparison use $o(1)$ instead of integrability for the low
top coefficients.  Products of finitely many lower-order norms and polynomial
losses are $e^{o(T')}$.  The parameter and pure target jets in spatial
source order $s-1$ have this property.  The smoothing, implicit position map
and compensator may require acceleration jets through $a^{(s+3)}$, accounting
for spatial derivatives of shifts and the two-order wave expression.  The
filter has worst differential order $1-P$; hence these jets require at most
\begin{equation}\label{co2:filter-count}
 \max(0,s+4-P)\le s-1
\end{equation}
derivatives of $\ell h$ once the \emph{single fixed} $P\ge5$ is chosen.
The larger fixed $P$ needed by the base resolvent tests only improves this
count.  Smooth probes, position integrations, exponential or bounded/reach
memories cost lower products and polynomials.  The original spatial equation
uses only low $c$ jets in its kernel families.

Solving for $\partial_t^2h$ in $H^{s-2}$ in a commutator gives
$C_s\sqrt{\mathcal E_s}+e^{o(T')}$, where $\mathcal E_s$ is the positive
first-jet energy.  Besides the differential pointwise solve, the order-$-1$
addition uses only its preceding induction order and coefficient products.
For example $K_{-1}\partial_t^2h$ needs $H^{s-3}$ on that jet and
coefficient algebra orders at most $s-2$.  Consequently, after exact
absorption of the differential principal part, the source difference at
$H^{s-1}$ costs
\[
 o(1)\sqrt{\mathcal E_s}+e^{o(T')}.
\]
Placements at orders already covered by induction give the additive term.
Top first derivatives multiply low vanishing differences.  Top second
derivatives in coefficient commutators or lower acceleration terms have the
one-order spatial gain and an $o(1)$ Lipschitz or coefficient difference.
Interior paraproduct pullbacks and differential currents also use only low
differences in front of top squared energy.  The outer products are estimated
in uniformly local unscaled coordinate patches of bounded overlap, and the
affine source factor changes density as in \eqref{co2:affine-density}.
Every top placement thus has a coefficient bounded, or vanishing, independently
of the terminal radius; high coefficients of order at most $s-1$ multiply
low arguments and belong to the additive induction bound.  If an interior
kernel or paraproduct acts on a low placement, the base estimates supply a
fixed compact Sobolev or H\"older buffer.  No polynomial in $T'$ multiplies
the top incremental coefficient.

It follows that $X_s=(1+\mathcal E_s)^{1/2}$ satisfies on the late interval
\begin{equation}\label{co2:higher-gronwall}
 \begin{gathered}
 \dot X_s\le\alpha_s(t)X_s+B_{s-1}(T'),
       \qquad B_{s-1}(T')=e^{o(T')},\\
 \alpha_s(t)\longrightarrow0
       \quad\text{independently of the terminal radius}.
 \end{gathered}
\end{equation}
For any prescribed $\sigma>0$, make $\alpha_s\le\sigma/3$ by a later
fixed starting time, use the induction with a rate smaller than $\sigma/3$
on $B_{s-1}$, and absorb the polynomial integration factor.  Gronwall gives
$X_s\le C_\sigma e^{\sigma T'}$.  This proves the spatial and first-time
orders.

Recover the higher time jets successively.  The normalized actual
$I+K_{-1}$ is uniformly invertible on spatial $L^2$ by low-order smallness.
At $H^{s-j}$ for $\partial_t^jh$, $2\le j\le s$, iterate its tame
negative-order gain; coefficient derivatives through at most
\[
 \max(s-j+1,4)\le s-1
\]
and fixed symbol seminorms suffice.  These have the induction bound.
The right side uses $c^{(j-1)}$ or lower, already recovered in earlier
time counts, while $p$ derivatives have the filter buffer
\eqref{co2:filter-count}.  Full mixed endpoints are therefore subexponential
without smallness of any high Sobolev norm of the gauge or datum.  Sobolev
embedding on enlarged regions and finite products give the stated smooth
bounds.  Every estimate was on a common finite slab; no limiting inner
horizon has entered the argument.
\end{proof}

The remaining task in the exterior analysis is to prove the far estimate and
close its coupling with \eqref{co2:sliced-acceleration} and
\eqref{co2:high-bound}.  The higher-order proposition will then apply to the
resulting solution in the same single base neighborhood.
